\section*{4 Using gate elimination to find functions in an optimal subcomputation}
Our approach to solving DecompProblem involves finding a ``test'' that functions in an optimal subcomputation pass but not too many other functions pass. The test will be based off the following function.
\setcounter{nttheorem}{21}
\begin{ntdefinition}[{Author Definition 22: $\Select[\cdot,\cdot]$}]
Let $f,g:\{0,1\}^n\to\{0,1\}$. We define the function $\Select[f,g]:\{0,1\}^n\times\{0,1\}\to\{0,1\}$ by
\[\Select[f,g](x,z)=\begin{cases}f(x),&\text{if }z=0,\\g(x),&\text{if }z=1.\end{cases}\]
\end{ntdefinition}
We emphasize that $\Select[f,g]$ function is only defined when $f$ and $g$ have the same arity.
Now, our ``test'' will be to see if the quantity $L(\Select[f,g])-L(f)$ is small (how small will depend on our reduction).
Indeed, for functions in an optimal subcomputation, this quantity is exactly one!\footnote{We will only prove that it is as most one, but the reader can check that if $g\ne f$ that a gate elimination argument actually implies equality.}
\begin{ntlemma}[Author Lemma 23]
Suppose $g$ is in an optimal subcomputation for $f$. Then $L(\Select[f,g])\leq L(f)+1$.
\end{ntlemma}
\begin{proof}
Since $g$ is in an optimal subcomputation for $f$, there exists an optimal formula $\phi=\phi_gO\phi_h$ such that $\phi_g$ computes $g$. If $O=\land$, then $\phi_g\land(\phi_h\lor z)$ is a formula for $\Select[f,g]$ of size $L(f)+1$. Otherwise $O=\lor$. Then $\phi_g\lor(\phi_h\land\neg z)$ is a formula for $\Select[f,g]$ of size $L(f)+1$.
\end{proof}
On the other hand, the number of functions that ``pass this test'' can be upper bounded in terms of $|\Canon{k}(f)|$.
\begin{ntlemma}[Author Lemma 24]
Let $k$ be a positive integer. Let $f:\{0,1\}^n\to\{0,1\}$. Assume $L(f)\geq 2$. Let $\operatorname{TestPassers}=\{g:L(\Select[f,g])-L(f)\leq k+1\}$. Then
\[|\operatorname{TestPassers}|\leq O(|\Canon{k}(f)|\cdot 2^kN\log N)\]
where $N=2^n$.
\end{ntlemma}
\begin{proof}
At a high-level the idea is that, given a formula $\phi$ of size $L(f)+k+1$ for computing $\Select[f,g]$, one can replace the $z$-leaves in $\phi$ with 0-leaves to obtain a formula $\phi'$ of size $L(f)+k+1$ for computing $f$ with at least one constant leaf. One can then use a careful gate elimination argument to remove precisely one constant leaf from $\phi'$ to obtain a formula $\phi''$ that still computes $f$ but has size $L(f)+k$. On the other hand, one can reverse this process by adding some constant leaf and gate to $\phi''$ and then replacing some subset of the constant leaves by $z$-leaves.
This gives us a way to describe any $g$ that passes the test, and thus allows us to bound the number of such $g$. In our bound, the $O(|\Canon{k}(f)|)$ factor corresponds to the choices for $\phi''$, the $O(N\log N)$ corresponds to the number of ways to add a new constant leaf and gate to $\phi''$ in order to obtain $\phi'$, and the $O(2^k)$ factor comes from the number of ways to chose a subset of the (at most $k$)0-leaves in $\phi'$ into $z$-leaves.
In detail, we prove this statement by giving a series of injections. At a high-level First, let $P$ denote the set of canonical formulas computing $f$ with size exactly $L(f)+k+1$ and with at least one constant-labelled leaf node in the formula.
We will give an injection from $\operatorname{TestPassers}$ to $P\times[2^{k+1}]$. Before defining our injection, we will need the following claims and definitions.
\begin{ntclaim}[Author Claim 25]
Suppose $\phi$ computes $\Select[f,g]$ and $f\ne g$, then $\phi$ has at least one leaf node labelled by $z$ or $\neg z$.
\end{ntclaim}
\begin{proof}
If $\phi$ does not have any $\{z,\neg z\}$ labelled leaves, then the output of $\phi$ does not depend on $z$. But this contradicts that $\phi$ computes $\Select[f,g]$ since $\Select[f,g]$ does depend on the $z$ input because $f\ne g$.
\end{proof}
\begin{ntclaim}[Author Claim 26]
Suppose $\phi\in P$. Then $\phi$ has at most $(k+1)$-many leaf nodes labelled by constants $\{0,1\}$.
\end{ntclaim}
\begin{proof}
Since $\phi\in P$, we know that $|\phi|=L(f)+k+1$ and $\phi$ computes $f$. Since $L(f)\geq 2$, we know that $|\phi|\geq k+3$.
If $\phi$ had more than $(k+1)$-many constant labelled leaves, it follows by a standard gate elimination argument (note here it is important that $|\phi|\geq k+3$) that there is a $\phi'$ that computes the same function as $\phi$ such that
\[|\phi'|<|\phi|-(k+1)<L(f).\]
But then $\phi$ would be a formula of size less than $L(f)$ computing $f$ which is a contradiction.
\end{proof}
We will also need the following definitions. Given a formula $\phi$ that can take $z$-variables as input, we define $\Subz(\phi)$ to be the formula $\phi'$ given by replacing the $z$-labeled leaves in $\phi$ with 0-labels and replacing the $(\neg z)$-labeled leaves in $\phi$ with 1-labels.
We note that the $\Subz$ operation in some sense respects formula isomorphisms.
\begin{ntclaim}[Author Claim 27]
Let $\phi$ be a formula of size $s$ that takes a $z$-variable as input. Let $\sigma:[2s-1]\to[2s-1]$ be a permutation. Then
\[\sigma\circ\Subz(\phi)=\Subz\circ\sigma(\phi).\]
\end{ntclaim}
\begin{proof}
The proof is essentially just applying the definition to both sides and seeing that the resulting edge sets and labelling functions are equal.
\end{proof}
We can also define a reverse operation to $\Subz$ as follows. Given a formula $\phi$ and a subset $S$ of leaf nodes in $\phi$ that are labelled by constants $\{0,1\}$, define $\Unsubz(\phi',S)$ to be the formula $\phi$ given by replacing 0-labeled leaves in $S$ with $z$-labels leaves and by replacing 1-labeled leaves in $S$ with $(\neg z)$-labels.
Indeed, the following claim whose proof we omit is easy to see.
\begin{ntclaim}[Author Claim 28]
For all formulas $\phi$ there exists a set $S$ such that
\[\phi=\Unsubz(\Subz(\phi),S).\]
\end{ntclaim}
Being more precise, $S$ is a subset of the leaf nodes in $\Subz(\phi)$ that are labelled by constants.
Now we are ready to describe our injection from $\operatorname{TestPassers}\to P\times[2^{k+1}]$ on an input $g\in\operatorname{TestPassers}$. Since $g\in\operatorname{TestPassers}$, there is a $\phi$ of size $L(f)+k+1$ computing $\Select[f,g]$. Let $\phi'=\Subz(\phi)$. Clearly $\phi'$ computes $f$ since $\phi$ computes $\Select[f,g]$. Let $\overline{\phi'}$ denote the canonical formula isomorphic to $\phi'$. Then there exists a permutation $\sigma$ such that
\begin{align*}\overline{\phi'}&=\sigma(\phi')\\&=\sigma\circ\Subz(\phi)\\&=\Subz\circ\sigma(\phi)\end{align*}
where the last equality comes from Claim 27. Thus, using Claim 28, we know that there exists a subset $S$ of the leaf nodes of $\overline{\phi'}$ labelled by constants such that
\[\Unsubz(\overline{\phi'},S)=\sigma(\phi).\]
Moreover, the set $S$ can be viewed as an element of $[2^{k+1}]$ because $\overline{\phi'}$ has at most $k+1$ leaf nodes labelled by constants. In particular, by construction, we have that
\[|\overline{\phi'}|=|\phi'|=|\phi|=L(f)+k+1,\]
and that $\overline{\phi'}$ computes $f$, so Claim 26 ensures that $\overline{\phi'}$ has at most $k+1$ many leaf nodes labelled by constants.
Hence, we define the output of our injection from $\operatorname{TestPassers}$ to $P\times[2^{k+1}]$ on input $g\in\operatorname{TestPassers}$ to be $(\overline{\phi'},S)$.
We must prove that this is indeed an injection. Towards this end, we claim that $\Unsubz(\overline{\phi'},S)$ is a formula computing $\Select[f,g]$. From this claim it is easy to see that this must be an injection.
\begin{ntclaim}[Author Claim 29]
$\Unsubz(\overline{\phi'},S)$ is a formula computing $\Select[f,g]$.
\end{ntclaim}
\begin{proof}
$S$ was chosen so that $\Unsubz(\overline{\phi'},S)=\sigma(\phi)$. Thus, $\Unsubz(\overline{\phi'},S)$ computes the same function as $\sigma(\phi)$ which in turn computes the same function as $\phi$, which computes $\Select[f,g]$ as desired.
\end{proof}
From this injection, we get that $|\operatorname{TestPassers}|\leq 2^{k+1}\cdot|P|$.
